\documentclass[11pt]{article}
\usepackage{labsheet_common}

\title{}
\date{}

\begin{document}
\thispagestyle{fancy}

\labtitle{AI Training -- Lab 1}{Estimating Truck Cargo Space with AI}{A hands-on introduction to labeling data and training a machine learning model}
\nocodebadge

You will follow copy-paste steps in a terminal, edit one plain-text settings
file, and label some photos by eye. That's it -- the AI code itself is
already written for you.

\section*{What you'll learn}
\begin{itemize}
  \item Why AI needs ``labeled'' examples before it can learn anything
  \item That your labeling choices directly change how good the AI becomes
  \item How changing a few settings (no programming) changes a model's accuracy
  \item How to read a simple accuracy score and compare it with classmates
\end{itemize}

\section*{The big picture}
We have photos taken by two cameras bolted inside a delivery truck's cargo
box (one near the front, one near the rear). As the truck gets loaded during
the day, the pile of boxes grows. Your goal: build an AI model that looks at
a new photo and guesses \textbf{what percentage of the cargo space is used}
(0\% = empty, 100\% = completely full).

\begin{notebox}
Nobody has ever written down the ``correct'' fill percentage for these
photos -- there's no automatic way to measure it. So before any AI can
learn, a \emph{human} has to look at each photo and estimate it. That's your
job in Part 1. This is normal in real AI projects: most of the effort goes
into preparing good data, not writing code.
\end{notebox}

\section*{What's already done for you}
\begin{itemize}
  \item 42 photos (21 front-camera + 21 rear-camera) have been organized into \path|datasets/processed/images/|.
  \item A spreadsheet-like file, \texttt{datasets/processed/manifest.csv}, lists every photo with an empty column waiting for your label.
  \item All the AI code -- feature extraction, model training, scoring, plotting -- is already written in \texttt{scripts/}. You will only \emph{run} it, never edit it.
  \item The only file you'll edit is \texttt{configs/cargo\_config.yaml}, a plain-text settings file (not code).
\end{itemize}

\section*{How it all fits together}
\begin{center}
\resizebox{\textwidth}{!}{%
\begin{tikzpicture}[node distance=6mm]
  \node[autobox] (raw) {Raw truck photos};
  \node[autobox, right=of raw] (organize) {Organize photos\\\texttt{organize\_\allowbreak dataset.py}};
  \node[actionbox, right=of organize] (label) {Label\\\texttt{fill\_level}};
  \node[autobox, right=of label] (split) {Split train/\allowbreak val/test\\\texttt{make\_\allowbreak splits.py}};

  \node[actionbox, below=15mm of split] (config) {Edit\\\texttt{cargo\_\allowbreak config.yaml}};
  \node[autobox, left=of config] (train) {Train model\\\texttt{cargo\_\allowbreak train.py}};
  \node[autobox, left=of train] (evaluate) {Evaluate model\\\texttt{cargo\_\allowbreak evaluate.py}};
  \node[scorebox, left=of evaluate] (score) {Your score:\\Mean Absolute Error};

  \draw[flowarrow] (raw) -- (organize);
  \draw[flowarrow] (organize) -- (label);
  \draw[flowarrow] (label) -- (split);
  \draw[flowarrow] (split) -- (config);
  \draw[flowarrow] (config) -- (train);
  \draw[flowarrow] (train) -- (evaluate);
  \draw[flowarrow] (evaluate) -- (score);

  \draw[loopA] (score.south) to[out=-55, in=-125, looseness=0.9]
    node[looplabel, midway, below=1mm] {try different settings (Part 5)} (config.south);
\end{tikzpicture}%
}
\end{center}
\diagramlegend

The top row happens once (organize the photos, then label them). The bottom
row is the part you'll repeat over and over in Part 5, trying to get the
lowest possible score.

\section*{Part 0: One-time setup}
Open a terminal (on Mac: the \textbf{Terminal} app; on Windows: \textbf{Anaconda Prompt} or \textbf{PowerShell}) and run these commands one at a time. Copy each line exactly, press Enter, and wait for it to finish before the next one.

\begin{lstlisting}
cd AI-training
python3 -m venv .venv
source .venv/bin/activate      # Windows: .venv\Scripts\activate
pip install -r requirements.txt
\end{lstlisting}

If that last command finishes without red error text, you're ready. You
only need to do Part 0 once (but you must run the \texttt{source
.venv/bin/activate} line again every time you open a new terminal window).

\section*{Part 1: Label the photos (this is the main task)}
\begin{enumerate}
  \item Open \texttt{datasets/processed/manifest.csv} in Excel, Numbers, or Google Sheets.
  \item For each row, open the photo it points to (the \texttt{filepath} column, inside \texttt{datasets/processed/}) and decide how full the truck looks, using this scale:
\end{enumerate}

\begin{center}
\begin{tabular}{@{}ll@{}}
\toprule
\textbf{Word you type in \texttt{fill\_level}} & \textbf{Meaning} \\
\midrule
\texttt{empty}  & Floor is fully visible, little to no cargo \\
\texttt{low}    & Some boxes/bags, most of the floor still visible \\
\texttt{medium} & Cargo covers roughly half the floor / waist-height pile \\
\texttt{high}   & Cargo fills most of the visible space, well above waist-height \\
\texttt{full}   & Cargo fills the frame edge-to-edge, no empty space visible \\
\bottomrule
\end{tabular}
\end{center}

\begin{enumerate}[resume]
  \item Type one of those five words into the \texttt{fill\_level} column for \textbf{every} row, then save the file as CSV (keep the filename \texttt{manifest.csv}).
\end{enumerate}

Full details and tips are in \texttt{docs/LABELING\_GUIDE.md} -- read it
before you start if anything above is unclear.

\begin{tipbox}
Your model can only ever be as good as your labels. If you label
carelessly, your AI will learn the wrong lesson -- same as it would for a
human being taught with wrong answer keys.
\end{tipbox}

\section*{Part 2: Split the data}
The AI needs to practice on some photos (``training'') and then be tested on
photos it has never seen (``testing'') -- otherwise it could just memorize
the answers instead of actually learning. Run:

\begin{lstlisting}
python3 scripts/make_splits.py
\end{lstlisting}

This automatically sorts your labeled photos into training/validation/test
groups and fills in a \texttt{fill\_pct} number for each label (e.g.
\texttt{medium} $\rightarrow$ 50). If it prints an error about missing
labels, go back to Part 1 -- every row must have a \texttt{fill\_level}
before this step will work.

\section*{Part 3: Train your model}
Open \texttt{configs/cargo\_config.yaml} in any plain-text editor (Notepad,
TextEdit, VS Code -- anything that isn't Word). You'll see settings like
this:

\begin{lstlisting}
model_type: random_forest
n_estimators: 100
max_depth: 5
random_seed: 42
\end{lstlisting}

You don't need to understand the AI math -- just try different values (see
the Glossary below for what each one means in plain language). Then run:

\begin{lstlisting}
python3 scripts/cargo_train.py
\end{lstlisting}

This trains a model using your current settings and saves it. It only takes
a few seconds.

\section*{Part 4: Evaluate your model}
\begin{lstlisting}
python3 scripts/cargo_evaluate.py
\end{lstlisting}

This tests your model on photos it never trained on, and prints something
like:

\begin{lstlisting}
Mean Absolute Error: 18.3 percentage points (lower is better)
\end{lstlisting}

That number is your score: on average, how many percentage points off your
model's guesses were. \textbf{Lower is better.} It also saves a picture,
\texttt{results/figures/cargo\_eval.png}, plotting your model's guesses
against the true labels -- dots closer to the diagonal line are better
predictions.

\section*{Part 5: Try to beat your own score (compete with classmates!)}
Go back to \texttt{configs/cargo\_config.yaml}, change one setting, then
re-run Part 3 and Part 4. Try:
\begin{itemize}
  \item Switching \texttt{model\_type} between \texttt{random\_forest}, \texttt{knn}, and \texttt{linear\_regression}
  \item For \texttt{random\_forest}: try different \texttt{n\_estimators} (e.g. 20 vs 100 vs 300) or \texttt{max\_depth} (e.g. 2 vs 5 vs 15)
  \item For \texttt{knn}: try different \texttt{n\_neighbors} (e.g. 1, 3, 10)
  \item Changing \texttt{random\_seed} to a different number
\end{itemize}

Keep a scoreboard with your class:

\begin{center}
\begin{tabular}{@{}*{4}{>{\raggedright\arraybackslash}p{3.2cm}}@{}}
\toprule
\textbf{Your name} & \textbf{model\_type} & \textbf{key settings} & \textbf{MAE (lower=better)} \\
\midrule
 & & & \\[10pt]
 & & & \\[10pt]
 & & & \\[10pt]
\bottomrule
\end{tabular}
\end{center}

Whoever gets the lowest Mean Absolute Error wins -- but also ask: did the
person with the ``best'' model also label their photos the most carefully?

\section*{Glossary}
\begin{itemize}
  \item \textbf{Model} -- the AI ``recipe'' that turns a photo into a fill-percentage guess. Training builds it; evaluating tests it.
  \item \textbf{Training set} -- photos the model is allowed to learn from.
  \item \textbf{Test set} -- photos held back so we can honestly check how good the model is on photos it has never seen.
  \item \textbf{Hyperparameter} -- a setting you choose before training (like \texttt{n\_estimators}), as opposed to something the model learns on its own.
  \item \textbf{Random Forest} -- a model made of many simple ``decision trees'' that each vote on an answer; the votes are averaged. \texttt{n\_estimators} = number of trees, \texttt{max\_depth} = how many yes/no questions each tree is allowed to ask.
  \item \textbf{KNN (k-nearest neighbors)} -- a model that finds the \texttt{n\_neighbors} most visually similar training photos and averages their fill percentages.
  \item \textbf{Linear Regression} -- the simplest possible model: it fits a straight-line relationship between photo features and fill percentage.
  \item \textbf{Overfitting} -- when a model memorizes quirks of the training photos instead of learning the general pattern, so it does great on training data but poorly on new photos. Very deep trees or too few neighbors are common causes.
  \item \textbf{MAE (Mean Absolute Error)} -- the average size of the model's mistakes, in percentage points. An MAE of 10 means the model is typically off by about 10 percentage points.
\end{itemize}

\begin{warnbox}
\begin{itemize}
  \item \textbf{``command not found: python3''} -- make sure you're in the \texttt{AI-training} folder and completed Part 0.
  \item \textbf{``No such file or directory: configs/cargo\_config.yaml''} -- you're probably running the command from the wrong folder; run \texttt{cd AI-training} first.
  \item \textbf{``manifest.csv has unlabeled/unsplit rows''} -- go back to Part 1, some rows are still missing a \texttt{fill\_level}, or you haven't run \texttt{scripts/make\_splits.py} yet.
  \item \textbf{Score looks wildly different from a classmate's with the same settings} -- check whether you both labeled the photos the same way. Small labeling differences are the most common cause, especially with only 42 photos total.
  \item \textbf{Everything says ``command not found'' after closing/reopening the terminal} -- re-run the \texttt{source .venv/bin/activate} line from Part 0.
\end{itemize}
\end{warnbox}

\section*{Something to think about (discuss with the class)}
\begin{itemize}
  \item We only have 42 photos from a single truck on a single day. What could go wrong if this model were used on a different truck, or at night?
  \item Why might two students who both hand-labeled the exact same photos still get slightly different \texttt{fill\_level} values?
  \item What would you need to collect to make this model trustworthy enough for real use?
\end{itemize}

\end{document}
